% BHU PYQ -- Transcribed & Corrected
% Paper: CHEMN42 : Techniques in Chemistry (End-Term)
% Exam: B.Sc. (Hons) Semester IV Examination, 2025-26 (NEP)
% Department: Chemistry

\documentclass[11pt,a4paper]{article}
\usepackage[a4paper,top=2.5cm,bottom=2.5cm,left=2.8cm,right=2.8cm,headheight=14pt]{geometry}
\usepackage{amsmath,amssymb}
\usepackage[T1]{fontenc}
\usepackage[utf8]{inputenc}
\usepackage{lmodern,microtype,fancyhdr,enumitem,hyperref,lastpage,parskip}

\hypersetup{
    colorlinks=true,
    urlcolor=black,
    linkcolor=black,
    pdftitle={CHEMN42: Techniques in Chemistry (End-Term) -- BHU Sem IV 2025-26 (NEP)}
}

\pagestyle{fancy}
\fancyhf{}
\renewcommand{\headrulewidth}{0.4pt}
\renewcommand{\footrulewidth}{0.4pt}
\fancyhead[L]{\small   \textit{Banaras Hindu University}}
\fancyhead[R]{\small   \textit{CHEMN42: Techniques in Chemistry (2025-26)}}
\fancyfoot[C]{\small Page  \thepage\ of \pageref{LastPage}}


\newlist{parts}{enumerate}{2}
\setlist[parts,1]{label=(\alph*),leftmargin=*,itemsep=4pt,topsep=2pt}
\setlist[parts,2]{label=(\roman*),leftmargin=*,itemsep=2pt,topsep=2pt}

\newcommand{\pts}[1]{\hfill   \textit{[#1]}}


\begin{document}


\begin{center}
  {\large\bfseries BANARAS HINDU UNIVERSITY}\\[2pt]
  {\normalsize B.Sc. (Hons.) --- Semester IV (NEP) Examination, 2025-26}\\[6pt]
  \rule{0.8\linewidth}{0.5pt}\\[6pt]
  {\large\bfseries CHEMISTRY}\\[4pt]
  {\large\bfseries Paper Code: CHEMN42 : Techniques in Chemistry}\\[4pt]
  {\normalsize Time: 2.5 Hours\hfill Maximum Marks: 70}\\[2pt]
  {\small REF NO. Conf/Even/N/2025-26/10}
\end{center}

\rule{\linewidth}{0.4pt}

\smallskip



\noindent   \textit{Instructions: Answer   \textbf{five} questions including Question No.~1 which is compulsory. All questions carry equal marks (14 marks each).}

\bigskip




\noindent   \textbf{Question 1.}   \textit{(Compulsory, 14 Marks Total)}

\begin{parts}
  \item What is a standard solution? Write the difference between primary and secondary standard solutions with suitable examples.\pts{2}
  \item What is an indicator? Describe its different types with suitable examples.\pts{2}
  \item Sketch the various types of bending vibration in IR of $ \text{H}_2 \text{O}$.\pts{2}
  \item Why pH ranges from 0 to 14 only?\pts{2}
  \item Calculate the pH value of $10^{-9}$ mole $ \text{HCl}$.\pts{2}
  \item What is crystallisation? State its importance in purification.\pts{2}
  \item What is the difference between qualitative and quantitative analysis? Explain with examples.\pts{2}
\end{parts}

\medskip\rule{\linewidth}{0.2pt}\medskip




\noindent   \textbf{Question 2.}

\begin{parts}
  \item Describe Instrumentation of UV-Visible absorption spectroscopy and its applications.\pts{4}
  \item Explain the concepts of molarity, normality and molality. Derive their expressions and discuss their interrelationships with suitable numerical examples.\pts{5}
  \item Explain the principle of complexometric titration. Describe the role of EDTA, metal ion indicators and applications of this method.\pts{5}
\end{parts}

\medskip\rule{\linewidth}{0.2pt}\medskip




\noindent   \textbf{Question 3.}

\begin{parts}
  \item Discuss the process of crystallization and precipitation. Explain the factors affecting crystal formation, purity and yield, along with their applications in chemical purification.\pts{6}
  \item Discuss the concept of conjugate acid-base pairs.\pts{2}
  \item Explain the general concepts of acids and bases according to Arrhenius, Br\o nsted-Lowry and Lewis theories, with suitable examples.\pts{6}
\end{parts}

\medskip\rule{\linewidth}{0.2pt}\medskip




\noindent   \textbf{Question 4.}

\begin{parts}
  \item Define the following terms used in UV-Visible absorption spectroscopy, with suitable examples:
  
\begin{parts}
    \item Bathochromic shift
    \item Hypsochromic shift
    \item Hypochromic effect
    \item Hyperchromic effect
  \end{parts}\pts{6}
  \item Discuss the principles and procedures of acid-base titrations. Explain the titration curves for a strong acid vs. a strong base and a weak base vs. a strong acid.\pts{4}
  \item A $25\,\mu \text{L}$ serum sample was analyzed for glucose content and found to contain $26.7\,\mu \text{g}$. Calculate the concentration of glucose in $\mu \text{g/mL}$.\pts{4}
\end{parts}

\medskip\rule{\linewidth}{0.2pt}\medskip




\noindent   \textbf{Question 5.}

\begin{parts}
  \item Describe the principle of chromatography. Compare paper chromatography, thin-layer chromatography (TLC) and column chromatography in terms of method, advantages and applications.\pts{8}
  \item What is the Lambert-Beer law? Explain each term in the expression.\pts{2}
  \item Write a short note on chemical shift in NMR.\pts{4}
\end{parts}

\medskip\rule{\linewidth}{0.2pt}\medskip




\noindent   \textbf{Question 6.}

\begin{parts}
  \item Describe infrared (IR) spectroscopy. Explain different IR sources. How we prepare sample for analysis?\pts{5}
  \item Describe redox titrations, including their principles, types and applications.\pts{5}
  \item Write notes on pH and pH calculations.\pts{4}
\end{parts}

\medskip\rule{\linewidth}{0.2pt}\medskip




\noindent   \textbf{Question 7.}

\begin{parts}
  \item Explain the principles of distillation. Describe in detail steam distillation and vacuum distillation, including their working mechanisms, advantages and applications.\pts{8}
  \item Differentiate between the fingerprint region and the group frequency region in IR spectroscopy.\pts{3}
  \item Explain precipitation titrations in details.\pts{3}
\end{parts}

\vfill
\rule{\linewidth}{0.4pt}

\begin{center}\small
\href{https://bhu-exam.vercel.app/}{Click here to visit our website}\\[4pt]
\href{https://me-aryan.vercel.app/}{Click here to visit our contributor}
\end{center}

\end{document}
